\documentclass[10pt,a4paper]{amsart}
\usepackage[margin=19mm]{geometry}
\usepackage{amsmath,amssymb,mathtools}
\usepackage[scr=rsfs,cal=euler]{mathalfa}
\usepackage{xfrac}
\usepackage[unicode,hidelinks,bookmarks=false]{hyperref}
\newcommand{\ie}{i.e.\ }
\newcommand{\cf}{cf.\ }
\newcommand{\resp}{respectively\ }
\newcommand{\Subsection}{\subsection}
\providecommand{\nobreakdash}{\nobreak\hbox{-}}
\setlength{\emergencystretch}{2em}
\multlinegap0pt
\usepackage{bm}
\usepackage{bbm}
\usepackage{dsfont}
\usepackage{xspace}
\usepackage{cancel}
\usepackage{mathtools}
\usepackage{amsthm}
\makeatletter
\@ifundefined{thm}{\newtheorem{thm}{Theorem}}{}
\@ifundefined{prop}{\newtheorem{prop}[thm]{Proposition}}{}
\makeatother
\DeclareMathOperator{\End}{End}
\title[Gradient flow in the kernel learning problem]{Gradient flow in the kernel learning problem}
\author{Yang Li, Feng Ruan}
\date{Source: arXiv:2506.08550v1 (CC BY 4.0)}
\begin{document}
\maketitle

\noindent\emph{Source and licence.} Gradient flow in the kernel learning problem. Version arXiv:2506.08550v1 (CC BY 4.0), distributed under the Creative Commons Attribution 4.0 International licence, as recorded in the arXiv licence metadata for this exact version and shown on the abstract page https://arxiv.org/abs/2506.08550v1. Full rights notice: licenses/arxiv-2506.08550-CC-BY-4.0.txt.\par\smallskip

\noindent\textit{Editorial scope.} Counted unit: the Prop whose source label is \texttt{prop:Euclideanflow} in the article (source0.tex lines 910--920, source label \texttt{prop:Euclideanflow}), delivered together with the article's own complete original proof of it (source0.tex lines 923--956, one \texttt{proof} environment of 34 lines). No mathematics is written, repaired or re-derived: statement and proof are the article's own text line for line. Required context retained uncounted: eqn:gradientflow (lines 288--290); eqn:Euclideanflow (lines 317--327); eqn:gradientflow2 (lines 631--633). Every definition, display and label of those lines is retained unchanged. Omitted from this package and not redistributed: 1--287, 291--316, 328--630, 634--909, 921--922, 957--1981. Nothing inside a retained excerpt is omitted or altered: every retained body is delivered exactly as the article's lines are written, so no omission receipt of any kind accompanies this package. Citations: the retained excerpts carry no citation command at all, so no bibliography entry of the article is transcribed and no reference is redistributed. Reference adaptation: none was needed and none was made; every reference inside the delivered unit resolves inside it, so no unresolved reference or citation is produced. Printed slips kept literal: no printed word, symbol, hypothesis or number of the counted statement or of its proof was altered. Layout and class: the article's own class is \texttt{article} with packages including babel, inputenc, amsmath, graphicx, amssymb, amsthm, tikz-cd, mathrsfs, todonotes, enumitem, yfonts, dsfont, MnSymbol, slashed; the wrapper is the lane's pinned \texttt{amsart} editorial wrapper at 10pt on A4, which restates the theorem-like environments the retained text needs and adds the article's own macro definitions that the retained text needs (\texttt{\textbackslash End}), with extra wrapper packages (bm, bbm, dsfont, xspace, cancel, mathtools). The delivered numbering is the unit's own. Assets: no retained span contains a figure environment, a graphics-inclusion command, a PDF-page inclusion, a table or a TikZ picture; no image file is copied into this package, the package declares no asset dependency and no image file is redistributed with it. Licence: version arXiv:2506.08550v1 of this article is distributed under the Creative Commons Attribution 4.0 International licence, as recorded in the arXiv licence metadata for this exact version and shown on the abstract page https://arxiv.org/abs/2506.08550v1. A full rights notice is delivered at licenses/arxiv-2506.08550-CC-BY-4.0.txt. Characters: every retained excerpt is pure ASCII, so the delivered engine, XeLaTeX with its default Latin Modern text font, reports no missing character.
	\begin{equation}\label{eqn:gradientflow}
		\frac{d}{dt} \Sigma(t)= - \nabla^{\mathfrak{g}} \mathcal{J}(\Sigma;\lambda).
	\end{equation}

	\begin{thm}
		(See Proposition \ref{prop:Euclideanflow})  Under the map 
		\[
		\pi: \End(V)\to Sym^2_{\geq 0}, \quad U\mapsto \Sigma=U^T U,
		\]
		the \emph{Euclidean gradient flow} on $\End(V)$ for the functional $J(U,\lambda)$
		\begin{equation}\label{eqn:Euclideanflow}
			\frac{dU_a^i}{dt}= - \frac{1}{4} C_{ac} C^{ij} (DJ(U,\lambda))_j^c  =  - \frac{1}{4} C_{ac} C^{ij}  \frac{\partial J(U,\lambda)}{\partial U_j^c}  ,
		\end{equation}
		is a lift of the Riemannian gradient flow.
	\end{thm}

	\begin{equation}\label{eqn:gradientflow2}
		\frac{d\Sigma_{ab}}{dt}= - \frac{1}{2} (\Sigma_{ac} C_{bd} (D \mathcal{J})^{cd}   +  \Sigma_{bd} C_{ac}  (D \mathcal{J})^{cd}    )
	\end{equation}

	\begin{prop}\label{prop:Euclideanflow}
		Under the map
		\[
		\pi: \End(V) \to Sym^2_{\geq 0}, \quad U\mapsto \Sigma= U^T U,
		\]
		the flow on $\End{V}$ defined by  
		\begin{equation}\label{eqn:Uflow}
			\frac{dU_a^i}{dt}= - \frac{1}{2} C_{ac} U_d^i (D\mathcal{J})^{cd}
		\end{equation}
		lifts the flow (\ref{eqn:gradientflow}) on $Sym^2_{\geq 0}$. Moreover, the flow (\ref{eqn:Uflow}) is equivalent to the \emph{Euclidean gradient flow} (\ref{eqn:Euclideanflow}) for the functional $J(U,\lambda)$ on $\End{V}$.
	\end{prop}

	\begin{proof}
		We compute in the index notation, to make manifest the role of the inner product $|\cdot|^2_V= C_{ij}dx^idx^j$; the interested reader may work out the matrix notation proof.
		
		
		
		Since $\Sigma_{ab}= U_a^i C_{ij} U_b^j$, we have $\frac{d}{dt}\Sigma_{ab}= \frac{d U_a^i}{dt} C_{ij} U_b^j+ U^i_a C_{ij} \frac{dU_b^j}{dt}$. We can rewrite the Riemannian gradient flow (\ref{eqn:gradientflow2}) in terms of $U$:
		\[
		\frac{d U_a^i}{dt} C_{ij} U_b^j+ U^i_a C_{ij} \frac{dU_b^j}{dt}= - \frac{1}{2} (U_a^i C_{ij} U_c^j C_{bd} (D \mathcal{J})^{cd}   + U_b^j U_d^i C_{ij} C_{ac}  (D \mathcal{J})^{cd}    ).
		\]
		This is implied by the flow (\ref{eqn:Uflow}) on $U$. This shows that the flow (\ref{eqn:Uflow}) projects down to the Riemannian gradient flow (\ref{eqn:gradientflow2}) under $\pi: \End(V)\to Sym^2_{\geq 0}$. 
		
		
		
		
		Furthermore,  we compute the differential of $J(U,\lambda)= \mathcal{J}(\Sigma)$ by the Leibniz rule and the chain rule:
		\[
		\begin{split}
			DJ= \pi^* D\mathcal{J}= & (D\mathcal{J})^{cd} \pi^* d\Sigma_{cd} = (D\mathcal{J})^{cd} d( U_c^i C_{ij} U_d^j  )
			\\
			= &   (D\mathcal{J})^{cd} C_{ij} U_d^j  dU_c^i  + (D\mathcal{J})^{cd} U_c^i C_{ij} dU_d^j 
			\\
			= &  2(D\mathcal{J})^{cd} C_{ij} U_d^j  dU_c^i  .
		\end{split}
		\]
		The last line here uses the symmetry of $(D\mathcal{J})^{cd}$ and $C_{ij}$. But $DJ= (DJ)_i^c dU_c^i$, so
		\[
		(DJ)_i^c = 2(D\mathcal{J})^{cd} C_{ij} U_d^j  ,\quad    C^{ij}(DJ)^c_j= 2(D\mathcal{J})^{cd} U_d^i. 
		\]
		Thus the flow (\ref{eqn:Uflow}) is equivalent to
		\[
		\frac{dU_a^i}{dt}= - \frac{1}{2} C_{ac} U_d^i (D\mathcal{J})^{cd}= - \frac{1}{4} C_{ac} C^{ij} (DJ)^c_j
		\]
		Up to a constant factor, this is just the Euclidean gradient flow for $J(U,\lambda)$ on the vector space $\End(V)$ with the Euclidean metric induced by $|\cdot|_V^2$.
	\end{proof}

\end{document}
